% ===========================================================================
%  Chapitre 14 — Variables aléatoires et lois de probabilité
%  PE 2026, composante TRAITEMENT DES DONNÉES — série OSE
% ===========================================================================
\bmtchapitre{Variables aléatoires et lois de probabilité}
  {Définir une variable aléatoire et sa loi, calculer espérance, variance et
   écart-type, reconnaître une épreuve de Bernoulli et choisir la loi
   binomiale lorsqu'elle s'applique.}
  {Traitement des données \textbullet\ environ 12 heures}

Le chapitre précédent mesurait des événements ; celui-ci mesure des
\emph{nombres}. À chaque résultat d'une expérience on associe une valeur
numérique — un gain, un effectif, une durée — et l'on étudie la façon dont
cette valeur se répartit. C'est ce qu'on appelle une variable aléatoire, et
c'est l'outil qui relie les probabilités à la statistique — le même outil
qu'utilise un assureur pour fixer une prime, ou une banque pour évaluer un
risque de défaut.

Deux nombres résument l'essentiel : l'espérance, qui donne la valeur moyenne
attendue sur un grand nombre de répétitions, et l'écart-type, qui mesure la
dispersion autour de cette moyenne. Une seule loi est au programme, la loi
binomiale, mais elle décrit une situation si fréquente — répéter la même
épreuve et compter les succès — qu'elle clôt presque tous les exercices de
probabilités du baccalauréat.

% ---------------------------------------------------------------------------
\begin{revision}
Chapitre précédent et statistique de première.

\begin{enumerate}[label=\textbf{\arabic*.}, leftmargin=*, itemsep=0.35em]
  \item Rappeler la définition de la moyenne d'une série statistique
        pondérée.
  \item Calculer $\combi{2}{5}$ et $\combi{3}{5}$.
  \item Que vaut $\proba{\evtc{A}}$ si $\proba{A} = \frac13$ ?
  \item Une urne contient trois boules blanches et sept noires. On en tire
        deux sans remise : probabilité d'obtenir deux noires ?
  \item Développer $(a+b)^{3}$ et repérer les coefficients binomiaux.
\end{enumerate}
\end{revision}

\begin{automatismes}
Sans calculatrice, sans brouillon.

\begin{enumerate}[label=\textbf{\alph*)}, leftmargin=*, itemsep=0.25em]
  \item $\combi{0}{4}$, $\combi{1}{4}$, $\combi{2}{4}$
  \item Somme des probabilités d'une loi
  \item $\Esp(X)$ si $X$ vaut $0$ ou $1$ avec probabilité $\frac12$
  \item $\Esp(X)$ pour $X \sim \loibin{10}{0{,}3}$
  \item $\Var(X)$ pour $X \sim \loibin{10}{0{,}3}$
  \item $\proba{X = 0}$ pour $X \sim \loibin{5}{\frac12}$
  \item $\Esp(2X+3)$ si $\Esp(X) = 4$
  \item $\ecart$ si $\Var(X) = 9$
\end{enumerate}
\end{automatismes}

\begin{corrige}[automatismes]
a) $1$, $4$, $6$ \quad b) $1$ \quad c) $\frac12$ \quad d) $3$ \quad
e) $2{,}1$ \quad f) $\frac{1}{32}$ \quad g) $11$ \quad h) $3$.
\end{corrige}

% ===========================================================================
\section{Définition d'une variable aléatoire}

\begin{definition}[variable aléatoire]
Une \textbf{variable aléatoire} sur un univers fini $\Omega$ est une
application $X$ de $\Omega$ dans $\R$. L'ensemble de ses valeurs se note
$X(\Omega)$.

L'événement « $X$ prend la valeur $x$ » se note $(X = x)$ ; c'est l'ensemble
des éventualités dont l'image par $X$ vaut $x$.
\end{definition}

\begin{avertissement}
Une variable aléatoire n'est ni une variable au sens de l'algèbre, ni un
nombre aléatoire : c'est une \emph{fonction}. Ce qui est aléatoire, c'est le
résultat de l'expérience ; $X$, elle, est parfaitement déterminée.
\end{avertissement}

% ===========================================================================
\section{Loi de probabilité et fonction de répartition}

\begin{definition}[loi de probabilité]
La \textbf{loi} de $X$ est la donnée, pour chaque valeur $x_{i}$ de
$X(\Omega)$, de la probabilité $p_{i} = \proba{X = x_{i}}$. On la présente
sous forme de tableau, et l'on a toujours
$\sum p_{i} = 1$.
\end{definition}

\begin{definition}[fonction de répartition]
La \textbf{fonction de répartition} de $X$ est
$F(x) = \proba{X \leq x}$, définie pour tout réel $x$. C'est une fonction en
escalier, croissante, nulle avant la plus petite valeur et égale à $1$ après
la plus grande.
\end{definition}

\begin{visualisation}[la loi et sa fonction de répartition]
\begin{center}
\begin{tikzpicture}
\begin{axis}[
  width = 6.1cm, height = 4.4cm,
  ybar, bar width = 12pt,
  axis lines = left,
  xlabel = {$x$}, ylabel = {$\proba{X = x}$},
  xmin = -0.7, xmax = 3.7, ymin = 0, ymax = 0.5,
  xtick = {0,1,2,3}, ytick = {0, 0.25, 0.5},
  axis line style = {bmtGris},
  label style = {font=\footnotesize, color=bmtGris},
  tick label style = {font=\footnotesize, color=bmtGris},
]
  \addplot[fill=bmtRaphia!45, draw=bmtRaphia] coordinates
    {(0,0.125) (1,0.375) (2,0.375) (3,0.125)};
\end{axis}
\begin{scope}[xshift=6.9cm]
\begin{axis}[
  width = 6.1cm, height = 4.4cm,
  axis lines = left,
  xlabel = {$x$}, ylabel = {$F(x)$},
  xmin = -0.7, xmax = 4.2, ymin = 0, ymax = 1.15,
  xtick = {0,1,2,3}, ytick = {0, 0.5, 1},
  axis line style = {bmtGris},
  label style = {font=\footnotesize, color=bmtGris},
  tick label style = {font=\footnotesize, color=bmtGris},
]
  \addplot[bmtIndigo, line width=1.2pt] coordinates {(-0.7,0) (0,0)};
  \addplot[bmtIndigo, line width=1.2pt] coordinates {(0,0.125) (1,0.125)};
  \addplot[bmtIndigo, line width=1.2pt] coordinates {(1,0.5) (2,0.5)};
  \addplot[bmtIndigo, line width=1.2pt] coordinates {(2,0.875) (3,0.875)};
  \addplot[bmtIndigo, line width=1.2pt] coordinates {(3,1) (4.1,1)};
  \addplot[bmtGris, dotted, line width=0.7pt] coordinates {(0,0) (0,0.125)};
  \addplot[bmtGris, dotted, line width=0.7pt] coordinates {(1,0.125) (1,0.5)};
  \addplot[bmtGris, dotted, line width=0.7pt] coordinates {(2,0.5) (2,0.875)};
  \addplot[bmtGris, dotted, line width=0.7pt] coordinates {(3,0.875) (3,1)};
\end{axis}
\end{scope}
\end{tikzpicture}
\end{center}

À gauche, la loi : une barre par valeur, de hauteur sa probabilité, et la
somme des hauteurs vaut $1$. À droite, la fonction de répartition, qui cumule
ces hauteurs de gauche à droite. Chaque saut a exactement la hauteur de la
barre correspondante — c'est la même information, lue de deux façons.
\end{visualisation}

\begin{exemple}[une loi complète]
Une urne contient trois boules blanches et sept noires. On tire deux boules
successivement et sans remise, et $X$ compte les blanches obtenues.
\[
  \proba{X = 0} = \frac{7}{10} \times \frac69 = \frac{7}{15},
  \qquad
  \proba{X = 1} = 2 \times \frac{3}{10} \times \frac79 = \frac{7}{15},
\]
\[
  \proba{X = 2} = \frac{3}{10} \times \frac29 = \frac{1}{15}.
\]
La somme vaut $\frac{7+7+1}{15} = 1$ : la loi est complète, et cette
vérification doit être écrite.
\end{exemple}

% ===========================================================================
\section{Espérance mathématique}

\begin{definition}[espérance]
\[
  \Esp(X) = \sum_{i} x_{i}\,p_{i} .
\]
\end{definition}

\begin{remarque}
L'espérance est la moyenne des valeurs, pondérée par leurs probabilités. Elle
s'interprète comme la valeur moyenne obtenue si l'on répétait l'expérience un
très grand nombre de fois. Elle n'est pas nécessairement une valeur possible
de $X$ : une espérance de $0{,}6$ enfant par famille ne choque personne.
\end{remarque}

\begin{propriete}[linéarité]
Pour tous réels $a$ et $b$, $\Esp(aX+b) = a\,\Esp(X)+b$.
\end{propriete}

\begin{methode}{décider si un jeu est équitable}
On modélise le gain algébrique du joueur par une variable aléatoire $G$ —
gains comptés positivement, mise comptée négativement — puis on calcule
$\Esp(G)$. Le jeu est équitable si $\Esp(G) = 0$, favorable au joueur si
$\Esp(G)>0$, défavorable sinon.
\end{methode}

% ===========================================================================
\section{Variance et écart-type}

\begin{definition}[variance, écart-type]
\[
  \Var(X) = \sum_{i}\left(x_{i}-\Esp(X)\right)^{2}p_{i}
          = \sum_{i}x_{i}^{2}p_{i}-\Esp(X)^{2},
  \qquad
  \ecart(X) = \sqrt{\Var(X)} .
\]
\end{definition}

\begin{propriete}[effet d'une transformation affine]
$\Var(aX+b) = a^{2}\Var(X)$ et $\ecart(aX+b) = \abs a\,\ecart(X)$ : ajouter
une constante ne change pas la dispersion.
\end{propriete}

\begin{avertissement}
La seconde expression de la variance — moyenne des carrés moins carré de la
moyenne — est la plus rapide, mais elle exige de la rigueur : on calcule
$\sum x_{i}^{2}p_{i}$, et non $\left(\sum x_{i}p_{i}\right)^{2}$. Et la
variance ne peut jamais être négative : un résultat négatif signale une erreur
de calcul, à corriger avant de poursuivre.
\end{avertissement}

\begin{entrainement}[loi, espérance, variance]
\exo[\niveau{1}] $X$ prend les valeurs $1$, $2$, $3$ avec les probabilités
$\frac12$, $\frac13$, $\frac16$. Calculer $\Esp(X)$ et $\Var(X)$.

\exo[\niveau{2}] Une urne contient neuf boules : quatre portent le numéro $1$,
trois le numéro $2$, deux le numéro $3$. On tire trois boules simultanément et
$X$ compte les boules numérotées $3$. Donner la loi de $X$ et son espérance.

\exo[\niveau{2}] Un jeu coûte $200$ ariary. On lance un dé : on gagne $600$
ariary si l'on obtient un $6$, rien sinon. Le jeu est-il équitable ?

\exo[\niveau{3}] Montrer que
$\Var(X) = \Esp\!\left(X^{2}\right)-\Esp(X)^{2}$.
\end{entrainement}

% ===========================================================================
\section{Épreuve et schéma de Bernoulli}

\begin{definition}[épreuve de Bernoulli]
Une \textbf{épreuve de Bernoulli} est une expérience aléatoire à deux issues,
appelées \textbf{succès} — de probabilité $p$ — et \textbf{échec}, de
probabilité $1-p$.

Un \textbf{schéma de Bernoulli} d'ordre $n$ est la répétition de $n$ épreuves
de Bernoulli identiques et \emph{indépendantes}.
\end{definition}

\begin{avertissement}
Les trois conditions sont indissociables : deux issues, même probabilité de
succès à chaque répétition, indépendance. Un tirage \emph{sans remise} viole
la troisième et parfois la deuxième : la loi binomiale ne s'y applique pas.
C'est le piège le plus fréquent de l'exercice.
\end{avertissement}

% ===========================================================================
\section{La loi binomiale}

\begin{theoreme}[loi binomiale]
Dans un schéma de Bernoulli d'ordre $n$ et de paramètre $p$, la variable
aléatoire $X$ comptant le nombre de succès suit la \textbf{loi binomiale}
$\loibin{n}{p}$ :
\[
  \proba{X = k} = \combi{k}{n}\,p^{k}(1-p)^{n-k},
  \qquad k \in \{0,1,\dots,n\}.
\]
De plus
\[
  \Esp(X) = np, \qquad \Var(X) = np(1-p), \qquad
  \ecart(X) = \sqrt{np(1-p)} .
\]
\end{theoreme}

\begin{demonstration}
Un chemin de l'arbre comportant exactement $k$ succès a pour probabilité
$p^{k}(1-p)^{n-k}$, puisque les épreuves sont indépendantes. Le nombre de tels
chemins est le nombre de façons de choisir les $k$ rangs des succès parmi $n$,
soit $\combi{k}{n}$. Ces chemins étant deux à deux incompatibles, on somme
leurs probabilités.
\end{demonstration}

\begin{methode}{reconnaître une situation binomiale}
Quatre questions, et il faut quatre « oui ».
\begin{enumerate}[label=\textbf{\arabic*.}, leftmargin=*, itemsep=0.15em]
  \item L'expérience se répète-t-elle un nombre fixé $n$ de fois ?
  \item Chaque répétition n'a-t-elle que deux issues ?
  \item La probabilité de succès est-elle la même à chaque fois ?
  \item Les répétitions sont-elles indépendantes ?
\end{enumerate}
Si oui, on écrit « $X$ suit la loi binomiale de paramètres $n$ et $p$ » — et
cette phrase, écrite explicitement, vaut des points.
\end{methode}

\begin{exemple}[une loi binomiale complète]
Une urne contient dix boules portant les lettres $A$, $A$, $A$, $B$, $B$,
$B$, $B$, $B$, $C$, $C$. On tire une boule, on la remet, et l'on recommence
cinq fois. Soit $X$ le nombre d'apparitions de $B$.

Les quatre conditions sont réunies, avec $n = 5$ et $p = \frac12$ :
$X \sim \loibin{5}{\frac12}$. Alors
\[
  \proba{X = 2} = \combi{2}{5}\left(\frac12\right)^{2}
  \left(\frac12\right)^{3} = \frac{10}{32} = \frac{5}{16},
\]
et $\proba{X \geq 1} = 1-\proba{X = 0} = 1-\frac{1}{32} = \frac{31}{32}$.
Enfin $\Esp(X) = 2{,}5$ et $\ecart(X) = \sqrt{5 \times \frac14}
= \frac{\sqrt5}{2}$.
\end{exemple}

\begin{entrainement}[loi binomiale]
\exo[\niveau{1}] $X \sim \loibin{6}{\frac13}$. Calculer $\proba{X = 2}$,
$\Esp(X)$ et $\Var(X)$.

\exo[\niveau{2}] Une machine produit $4\,\%$ de pièces défectueuses. On en
prélève $20$. Probabilité d'en trouver exactement une défectueuse ? Au moins
une ?

\exo[\niveau{2}] $X \sim \loibin{n}{0{,}05}$. Déterminer le plus petit $n$ tel
que $\proba{X \geq 1} \geq 0{,}99$.

\exo[\niveau{3}] $X \sim \loibin{n}{p}$. Démontrer la formule
$\Esp(X) = np$ dans le cas $n = 2$, puis expliquer comment la récurrence
s'établit.
\end{entrainement}

% ===========================================================================
\section{Choisir la loi adaptée}

\begin{propriete}[deux situations, deux traitements]
\begin{center}
\setlength{\tabcolsep}{8pt}
\renewcommand{\arraystretch}{1.3}
\begin{tabular}{@{}lll@{}}
\toprule
\textsf{\footnotesize\bfseries SITUATION} &
\textsf{\footnotesize\bfseries INDÉPENDANCE} &
\textsf{\footnotesize\bfseries TRAITEMENT} \\
\midrule
Tirages avec remise & oui & loi binomiale \\
Tirages sans remise & non & dénombrement direct \\
Tirage simultané & non & dénombrement direct \\
\bottomrule
\end{tabular}
\end{center}
\end{propriete}

\begin{remarque}
Lorsque la population est très grande devant le nombre de prélèvements, un
tirage sans remise se comporte presque comme un tirage avec remise : l'énoncé
autorise alors explicitement l'assimilation à un schéma de Bernoulli. Cette
autorisation doit figurer dans le texte ; on ne se l'accorde pas soi-même.
\end{remarque}

% ===========================================================================
\begin{annales}
La variable aléatoire clôt l'exercice de probabilités : on demande sa loi,
puis son espérance ou son écart-type, et parfois une inéquation sur le nombre
de répétitions. Les énoncés qui suivent sont repris de sujets réellement
posés.

\exoann{Baccalauréat, Madagascar, série S, session 2021 — exercice 1, partie II}
Quatre billes numérotées $1$ à $4$ sont lancées vers quatre trous $T_{1}$,
$T_{2}$, $T_{3}$, $T_{4}$ ; chaque bille entre dans un trou et tous les
événements élémentaires sont équiprobables. Soit $X$ la variable aléatoire
égale au nombre de billes reçues par le trou $T_{3}$ lors d'un lancer. Donner
la loi de probabilité de $X$.

\exoann{Baccalauréat, Madagascar, série S, session 2023 — exercice 1, partie B}
\emph{On admet que l'épreuve $E$ du sujet réalise l'événement $A$ avec la
probabilité $\frac14$.}

On répète trois fois de suite et de manière indépendante l'épreuve $E$. Soit
$X$ la variable aléatoire égale au nombre de réalisations de $A$.
\begin{enumerate}[label=\textbf{\alph*)}, leftmargin=*, itemsep=0.15em]
  \item Donner la loi de $X$.
  \item Calculer l'écart-type de $X$.
\end{enumerate}

\exoann{Baccalauréat blanc, lycée Philibert Tsiranana, Terminale S, 2022-2023 — exercice 1, partie II}
Une urne contient quatre boules numérotées $1$, trois numérotées $2$ et deux
numérotées $3$. On tire simultanément trois boules et l'on note $X$ le nombre
de boules numérotées $3$. Donner la loi de probabilité de $X$.

\exoann{Baccalauréat, Congo-Brazzaville, série C, session 2023 — exercice 4}
Une urne contient dix boules, trois blanches et sept noires. On effectue deux
tirages successifs sans remise, et $X$ compte les boules blanches obtenues.
Après avoir donné la loi de $X$, calculer son espérance mathématique.

\exoann{Baccalauréat, Mauritanie, série C, session 2022 — exercice 1}
Cinq pour cent des candidats au baccalauréat d'une session donnée étaient en
série C. On choisit au hasard $n$ candidats ($n \geq 2$), les tirages étant
effectués avec remise, et l'on note $X$ le nombre de candidats de série C
obtenus.
\begin{enumerate}[label=\textbf{\alph*)}, leftmargin=*, itemsep=0.15em]
  \item Déterminer la loi de probabilité de $X$ et son espérance
        mathématique.
  \item Calculer $\proba{A}$ et $\proba{B}$, où $A$ est « exactement deux
        candidats sont en série C » et $B$ « au moins un candidat est en série
        C ».
\end{enumerate}

\exoann{Baccalauréat, Mauritanie, série C, session 2022 — exercice 1, suite}
On pose $p_{n} = \proba{B}$.
\begin{enumerate}[label=\textbf{\alph*)}, leftmargin=*, itemsep=0.15em]
  \item Calculer $\lim_{n\to+\infty}p_{n}$ et interpréter le résultat.
  \item Quel nombre minimal de candidats faut-il choisir pour que
        $p_{n} \geq 0{,}99$ ?
\end{enumerate}

\exoann{Baccalauréat, Cameroun, séries D et TI, session 2024 — exercice 2}
Une loterie comporte vingt billets : un billet fait gagner $1\,000$~F, trois
billets font gagner $500$~F, les autres ne rapportent rien. Une personne
achète deux billets. Soit $X$ la variable aléatoire donnant la somme totale
gagnée.
\begin{enumerate}[label=\textbf{\alph*)}, leftmargin=*, itemsep=0.15em]
  \item Indiquer toutes les valeurs possibles de $X$.
  \item Dresser le tableau de la loi de probabilité de $X$.
  \item Calculer $\proba{X>500}$.
\end{enumerate}

\exoann{D'après le baccalauréat, Cameroun, séries C et E, session 2024 — exercice 1}
Une urne contient six boules numérotées de $1$ à $6$. On tire successivement
et avec remise deux boules, dont on note les numéros $a$ et $b$, puis on
considère l'équation caractéristique $r^{2}+ar+b = 0$. On admet que la
probabilité que cette équation n'ait pas de racine réelle vaut
$\frac{17}{36}$. Combien de répétitions indépendantes de cette expérience
garantissent au moins $98\,\%$ de chances d'obtenir au moins une fois deux
solutions non réelles ?
\end{annales}

\begin{corrige}[annales]
\textbf{1.} Chaque bille choisit son trou indépendamment des autres, avec la
probabilité $\frac14$ de tomber dans $T_{3}$ : $X$ suit la loi binomiale de
paramètres $4$ et $\frac14$. D'où
\[
  \proba{X = 0} = \frac{81}{256}, \quad
  \proba{X = 1} = \frac{108}{256}, \quad
  \proba{X = 2} = \frac{54}{256},
\]
\[
  \proba{X = 3} = \frac{12}{256}, \quad
  \proba{X = 4} = \frac{1}{256},
\]
et la somme vaut bien $1$.

\textbf{2. a)} $X \sim \loibin{3}{\frac14}$ :
$\proba{X = 0} = \frac{27}{64}$, $\proba{X = 1} = \frac{27}{64}$,
$\proba{X = 2} = \frac{9}{64}$, $\proba{X = 3} = \frac{1}{64}$.
\textbf{b)} $\Var(X) = 3 \times \frac14 \times \frac34 = \frac{9}{16}$, donc
$\ecart(X) = \frac34$.

\textbf{3.} Le tirage est simultané : on dénombre. Il y a
$\combi{3}{9} = 84$ tirages possibles, et
\[
  \proba{X = 0} = \frac{\combi{3}{7}}{84} = \frac{35}{84} = \frac{5}{12},
  \quad
  \proba{X = 1} = \frac{2\combi{2}{7}}{84} = \frac12,
  \quad
  \proba{X = 2} = \frac{7}{84} = \frac{1}{12}.
\]
L'espérance vaut $\frac12+\frac{2}{12} = \frac23$.

\textbf{4.} $\proba{X = 0} = \frac{7}{15}$, $\proba{X = 1} = \frac{7}{15}$,
$\proba{X = 2} = \frac{1}{15}$, donc
$\Esp(X) = \frac{7}{15}+\frac{2}{15} = \frac{9}{15} = \frac35$.

\textbf{5. a)} Les tirages étant effectués avec remise, les épreuves sont
identiques et indépendantes : $X \sim \loibin{n}{0{,}05}$ et
$\Esp(X) = 0{,}05\,n$.
\textbf{b)} $\proba{A} = \combi{2}{n}(0{,}05)^{2}(0{,}95)^{n-2}$ et
$\proba{B} = 1-(0{,}95)^{n}$.

\textbf{6. a)} $(0{,}95)^{n} \to 0$, donc $p_{n} \to 1$ : en interrogeant
suffisamment de candidats, on est pratiquement certain d'en rencontrer au
moins un de série C.
\textbf{b)} $(0{,}95)^{n} \leq 0{,}01$ donne
$n \geq \frac{\ln0{,}01}{\ln0{,}95} \approx 89{,}8$ : il faut au minimum
$90$ candidats.

\textbf{7. a)} Les deux billets sont pris parmi vingt ; les gains possibles
sont $0$, $500$, $1\,000$ et $1\,500$~F — la valeur $2\,000$ est impossible,
un seul billet valant $1\,000$~F.
\textbf{b)} Il y a $\combi{2}{20} = 190$ achats possibles :
\[
  \proba{X = 0} = \frac{12}{19}, \quad
  \proba{X = 500} = \frac{24}{95}, \quad
  \proba{X = 1000} = \frac{1}{10}, \quad
  \proba{X = 1500} = \frac{3}{190}.
\]
\textbf{c)} $\proba{X>500} = \frac{19+3}{190} = \frac{22}{190}
= \frac{11}{95}$.

\textbf{8.} Chaque expérience est une épreuve de Bernoulli de paramètre
$p = \frac{17}{36}$, et les répétitions sont indépendantes. La probabilité de
n'obtenir aucun succès en $n$ répétitions vaut
$\left(\frac{19}{36}\right)^{n}$, et l'on veut
\[
  1-\left(\frac{19}{36}\right)^{n} \geq 0{,}98
  \iff \left(\frac{19}{36}\right)^{n} \leq 0{,}02,
\]
c'est-à-dire $n \geq \dfrac{\ln 0{,}02}{\ln\frac{19}{36}} \approx 6{,}12$.
Sept répétitions suffisent.
\end{corrige}

% ===========================================================================
\begin{jecherche}[la tontine du quartier]
Douze familles d'un quartier cotisent chaque mois à une tontine. Chaque mois,
une famille est tirée au sort pour recevoir la cagnotte ; le tirage se fait
avec remise, une même famille pouvant donc être servie plusieurs fois dans
l'année.

\begin{enumerate}[label=\textbf{\arabic*.}, leftmargin=*, itemsep=0.3em]
  \item Soit $X$ le nombre de fois où une famille donnée est servie au cours
        des douze mois. Justifier que $X$ suit une loi binomiale et préciser
        ses paramètres.
  \item Calculer l'espérance et l'écart-type de $X$.
  \item Calculer la probabilité qu'une famille donnée ne soit jamais servie
        dans l'année.
  \item Calculer la probabilité qu'elle soit servie exactement deux fois.
  \item Les familles proposent de tirer sans remise, chacune étant alors
        servie une fois et une seule. Comparer les deux systèmes du point de
        vue de l'espérance, puis de la dispersion.
  \item Combien de mois faudrait-il, avec remise, pour qu'une famille donnée
        ait plus de $90\,\%$ de chances d'être servie au moins une fois ?
\end{enumerate}
\end{jecherche}

\begin{corrige}[la tontine du quartier]
\textbf{1.} Chaque mois est une épreuve à deux issues — la famille est servie
ou non — de probabilité constante $\frac{1}{12}$, et les tirages avec remise
sont indépendants. Donc $X \sim \loibin{12}{\frac{1}{12}}$.

\textbf{2.} $\Esp(X) = 12 \times \frac{1}{12} = 1$ : en moyenne, une famille
est servie une fois par an. La variance vaut
$12 \times \frac{1}{12} \times \frac{11}{12} = \frac{11}{12}$, donc
$\ecart(X) \approx 0{,}957$.

\textbf{3.} $\proba{X = 0} = \left(\frac{11}{12}\right)^{12} \approx 0{,}352$
— plus d'une chance sur trois de n'être jamais servi, alors même que
l'espérance vaut $1$.

\textbf{4.} $\proba{X = 2} = \combi{2}{12}\left(\frac{1}{12}\right)^{2}
\left(\frac{11}{12}\right)^{10} \approx 0{,}213$.

\textbf{5.} Sans remise, chaque famille est servie exactement une fois :
l'espérance vaut encore $1$, mais l'écart-type est nul. Les deux systèmes
sont équivalents en moyenne et radicalement différents en équité : c'est
précisément pourquoi les tontines réelles procèdent sans remise.

\textbf{6.} Il faut $1-\left(\frac{11}{12}\right)^{n} \geq 0{,}9$, soit
$\left(\frac{11}{12}\right)^{n} \leq 0{,}1$, c'est-à-dire
$n \geq \frac{\ln 0{,}1}{\ln\frac{11}{12}} \approx 26{,}5$ : il faudrait
vingt-sept mois, soit plus de deux ans.
\end{corrige}

% ===========================================================================
\begin{jemeteste}
Répondre par vrai ou faux, en justifiant en une ligne.

\begin{enumerate}[label=\textbf{\arabic*.}, leftmargin=*, itemsep=0.28em]
  \item L'espérance est toujours une valeur prise par la variable.
  \item La somme des probabilités d'une loi vaut $1$.
  \item Un tirage sans remise conduit à une loi binomiale.
  \item $\Var(X)$ peut être négative.
  \item Pour $X \sim \loibin{n}{p}$, $\Esp(X) = np$.
  \item $\Esp(3X-2) = 3\Esp(X)-2$.
\end{enumerate}
\end{jemeteste}

\begin{corrige}[je me teste]
\textbf{1.} Faux — elle peut ne correspondre à aucune valeur possible.
\textbf{2.} Vrai.
\textbf{3.} Faux — l'indépendance est perdue ; on dénombre directement.
\textbf{4.} Faux — c'est une somme de carrés pondérés par des probabilités.
\textbf{5.} Vrai.
\textbf{6.} Vrai, par linéarité de l'espérance.
\end{corrige}

% ===========================================================================
\begin{commeaubac}
\bmtsource{Baccalauréat de l'Enseignement Général, Madagascar, série OSE, option OSE, session 2026 — exercice 1}
Une urne contient trois plaquettes indiscernables au toucher, numérotées $1$,
$2$ et $3$. On tire successivement et avec remise trois plaquettes ; on note
dans l'ordre les trois numéros obtenus. On considère les événements
\[
  A : \text{« au moins une plaquette porte un numéro pair »},
\]
\[
  B : \text{« le produit des trois numéros obtenus vaut $6$ »} .
\]
Soit $X$ la variable aléatoire égale au produit des trois numéros obtenus.

\begin{enumerate}[label=\textbf{\arabic*.}, leftmargin=*, itemsep=0.25em]
  \item Calculer $\proba{A}$ et $\proba{B}$. \hfill\textup{(1 pt)}
  \item Déterminer $X(\Omega)$, puis dresser le tableau de la loi de
        probabilité de $X$. \hfill\textup{(1,25 pt)}
  \item Calculer $\Esp(X)$ et $\Var(X)$. \hfill\textup{(0,75 pt)}
  \item On note $F$ la fonction de répartition de $X$. Représenter
        graphiquement $F$. \hfill\textup{(0,5 pt)}
  \item On répète trois fois, de manière indépendante, l'épreuve consistant à
        tirer les trois plaquettes avec remise. Calculer la probabilité
        d'obtenir exactement deux fois l'événement $B$. \hfill\textup{(0,5 pt)}
\end{enumerate}
\end{commeaubac}

\begin{corrige}[comme au baccalauréat]
\textbf{1.} Chaque tirage est indépendant et uniforme sur $\{1\,;2\,;3\}$,
dont un seul numéro — le $2$ — est pair. L'événement contraire de $A$ est
« aucune des trois plaquettes ne porte le numéro $2$ », de probabilité
$\left(\frac23\right)^{3} = \frac{8}{27}$, donc
\[
  \proba{A} = 1-\frac{8}{27} = \frac{19}{27}.
\]
Le produit vaut $6$ exactement lorsque les trois numéros tirés sont $1$, $2$
et $3$ dans un ordre quelconque : il y a $3! = 6$ triplets ordonnés
favorables sur $3^{3} = 27$ triplets possibles, donc
$\proba{B} = \dfrac{6}{27} = \dfrac{2}{9}$.

\textbf{2.} En dénombrant les $27$ triplets ordonnés selon leur composition
non ordonnée, les produits possibles et leurs effectifs sont
\[
  \begin{array}{c|cccccccccc}
    x_{i} & 1 & 2 & 3 & 4 & 6 & 8 & 9 & 12 & 18 & 27 \\
    \hline
    \text{effectif} & 1 & 3 & 3 & 3 & 6 & 1 & 3 & 3 & 3 & 1
  \end{array}
\]
(par exemple $x = 6$ provient des six permutations de $(1,2,3)$, et
$x = 8$ du seul triplet $(2,2,2)$), d'où $X(\Omega)
= \{1,2,3,4,6,8,9,12,18,27\}$ et la loi
\[
  \begin{array}{c|cccccccccc}
    x_{i} & 1 & 2 & 3 & 4 & 6 & 8 & 9 & 12 & 18 & 27 \\
    \hline
    \proba{X = x_{i}} & \frac{1}{27} & \frac{3}{27} & \frac{3}{27} &
    \frac{3}{27} & \frac{6}{27} & \frac{1}{27} & \frac{3}{27} &
    \frac{3}{27} & \frac{3}{27} & \frac{1}{27}
  \end{array}
\]
La somme des effectifs vaut bien $1+3+3+3+6+1+3+3+3+1 = 27$.

\textbf{3.} Les trois tirages étant indépendants et identiquement
distribués, $\Esp(X) = \Esp(X_{1})^{3}$ où $X_{1}$ est le numéro d'un seul
tirage : $\Esp(X_{1}) = \frac{1+2+3}{3} = 2$, donc $\Esp(X) = 2^{3} = 8$. On
retrouve ce résultat directement à partir du tableau :
\[
  \Esp(X) = \frac{1 \times1+2\times3+3\times3+4\times3+6\times6+8\times1
  +9\times3+12\times3+18\times3+27\times1}{27} = \frac{216}{27} = 8 .
\]
De même $\Esp\!\left(X_{1}^{2}\right) = \frac{1+4+9}{3} = \frac{14}{3}$, et
par indépendance
$\Esp\!\left(X^{2}\right) = \left(\frac{14}{3}\right)^{3}
= \frac{2744}{27}$, d'où
\[
  \Var(X) = \Esp\!\left(X^{2}\right)-\Esp(X)^{2}
  = \frac{2744}{27}-64 = \frac{2744-1728}{27} = \frac{1016}{27}
  \approx 37{,}6 .
\]

\textbf{4.} En cumulant les probabilités de gauche à droite :
\begin{center}
\begin{tikzpicture}
\begin{axis}[
  width = 11.5cm, height = 5.2cm,
  axis lines = left,
  xlabel = {$x$}, ylabel = {$F(x)$},
  xmin = -1, xmax = 29, ymin = 0, ymax = 1.15,
  xtick = {1,2,3,4,6,8,9,12,18,27}, ytick = {0, 0.5, 1},
  axis line style = {bmtGris},
  label style = {font=\footnotesize, color=bmtGris},
  tick label style = {font=\scriptsize, color=bmtGris},
]
  \addplot[bmtIndigo, line width=1.1pt] coordinates {(-1,0) (1,0)};
  \addplot[bmtIndigo, line width=1.1pt] coordinates {(1,0.037) (2,0.037)};
  \addplot[bmtIndigo, line width=1.1pt] coordinates {(2,0.148) (3,0.148)};
  \addplot[bmtIndigo, line width=1.1pt] coordinates {(3,0.259) (4,0.259)};
  \addplot[bmtIndigo, line width=1.1pt] coordinates {(4,0.370) (6,0.370)};
  \addplot[bmtIndigo, line width=1.1pt] coordinates {(6,0.593) (8,0.593)};
  \addplot[bmtIndigo, line width=1.1pt] coordinates {(8,0.630) (9,0.630)};
  \addplot[bmtIndigo, line width=1.1pt] coordinates {(9,0.741) (12,0.741)};
  \addplot[bmtIndigo, line width=1.1pt] coordinates {(12,0.852) (18,0.852)};
  \addplot[bmtIndigo, line width=1.1pt] coordinates {(18,0.963) (27,0.963)};
  \addplot[bmtIndigo, line width=1.1pt] coordinates {(27,1) (29,1)};
  \addplot[bmtGris, dotted, line width=0.6pt] coordinates {(1,0) (1,0.037)};
  \addplot[bmtGris, dotted, line width=0.6pt] coordinates {(2,0.037) (2,0.148)};
  \addplot[bmtGris, dotted, line width=0.6pt] coordinates {(3,0.148) (3,0.259)};
  \addplot[bmtGris, dotted, line width=0.6pt] coordinates {(4,0.259) (4,0.370)};
  \addplot[bmtGris, dotted, line width=0.6pt] coordinates {(6,0.370) (6,0.593)};
  \addplot[bmtGris, dotted, line width=0.6pt] coordinates {(8,0.593) (8,0.630)};
  \addplot[bmtGris, dotted, line width=0.6pt] coordinates {(9,0.630) (9,0.741)};
  \addplot[bmtGris, dotted, line width=0.6pt] coordinates {(12,0.741) (12,0.852)};
  \addplot[bmtGris, dotted, line width=0.6pt] coordinates {(18,0.852) (18,0.963)};
  \addplot[bmtGris, dotted, line width=0.6pt] coordinates {(27,0.963) (27,1)};
\end{axis}
\end{tikzpicture}
\end{center}
$F$ est nulle avant $1$, vaut $1$ à partir de $27$, et effectue dix sauts,
chacun de hauteur la probabilité de la valeur franchie.

\textbf{5.} L'événement $B$ — « produit égal à $6$ » — a pour probabilité
$p = \frac29$ à chaque répétition de l'épreuve complète, et les trois
répétitions sont indépendantes : le nombre $Y$ de réalisations de $B$ suit
la loi $\loibin{3}{\frac29}$. Donc
\[
  \proba{Y = 2} = \combi{2}{3}\left(\frac29\right)^{2}\left(\frac79\right)
  = 3 \times \frac{4}{81} \times \frac79 .
\]
\[
  \proba{Y = 2} = \frac{84}{729} = \frac{28}{243} \approx 0{,}115 .
\]
\end{corrige}

% ===========================================================================
\begin{notepourlaclasse}
Le chapitre se joue sur une seule décision : la situation est-elle binomiale
ou non ? Faites-en un rituel de début d'exercice — les quatre questions,
posées à voix haute, avant tout calcul. Les élèves qui appliquent la formule
binomiale à un tirage simultané perdent l'exercice entier, et ils sont
nombreux.

L'écart-type est souvent calculé mécaniquement sans jamais être interprété.
Une séance consacrée à comparer deux lois de même espérance et d'écarts-types
très différents — la tontine avec et sans remise, par exemple — vaut mieux
qu'une démonstration supplémentaire.

Enfin, l'exercice corrigé du sujet 2026 mérite d'être retravaillé intégralement
en classe : il enchaîne exactement la structure attendue — probabilités
simples, loi complète d'une variable produit, espérance et variance,
fonction de répartition, puis un second niveau d'aléa par répétition de
l'épreuve. C'est la meilleure préparation directe à l'exercice 1 de
l'épreuve.
\end{notepourlaclasse}
